\begin{frame}{$\varepsilon$ 를 줄여가며 학습: 탐험에서 착취로}
  $\varepsilon = 0.1$ 을 고집해서 40만 에피소드까지 돌리면,
  학습은 ``대략 좋은 정책''에서 멈추기 쉽다 --- 매 스텝 10\%
  의 무작위 행동을 계속 섞어주니 $Q$ 추정치가 항상 조금씩
  흔들리고 \textbf{경계 상태의 판단이 굳어지지 않기} 때문이다.
  \vskip0.5em
  실전에서는 보통 \textbf{학습이 진행될수록 $\varepsilon$ 를
  줄인다} (\kb{$\varepsilon$-decay}, 탐험 스케줄) --- 초반에
  $\varepsilon$ 을 크게 해서 전략 공간을 넓게 훑고, 후반에는
  줄여 ``지금까지의 학습''을 믿고 행동.
  \vskip0.5em
  이 실습: 30만 에피소드($\varepsilon = 0.1$) 뒤에 10만
  에피소드를 $\varepsilon = 0.02$ 로 추가 학습. 효과(작지만
  실제로 측정됨):
  \begin{itemize}
    \item 정확 정책과의 일치율: \textbf{104/110 $\to$ 105/110}.
    \item $|Q - Q^{*}|$ 의 평균: 0.029 $\to$ 0.027.
  \end{itemize}
\end{frame}
